\clearpage
\section{Introduction}

This assignment extends the \textbf{Smart Home Alarm System} developed in
\textbf{assignment 3} and turns it into a \textbf{distributed typed-actor system}
based on \textbf{Apache Pekko Cluster}. The required alarm logic does not
change: the same commands, delays, and alarm transitions must still hold.

The central idea of the project is simple: \textbf{same alarm logic, different
execution model}. Instead of one local actor system, the implementation uses
multiple cluster nodes, each started with \textbf{one explicit role}. In the
final solution this leads to a clear deployment choice: \textbf{one node, one
main role}, so that control unit, keypad, and sensors can be started and
observed independently.

The most important architectural consequence is that the \textbf{Control Unit}
can no longer be just a normal local actor. Since the alarm state must still
have a \textbf{single logical owner}, the clustered version models it as a
\textbf{cluster singleton}. This keeps the state authority unique even if
multiple eligible nodes exist in the cluster.

The clustered version therefore focuses on three technical aspects:
\begin{itemize}
    \item \textbf{node startup}, because each JVM must join the same cluster with the correct role;
    \item \textbf{remote discovery}, because peripheral actors must find the control unit dynamically;
    \item \textbf{safe recovery}, because a recreated control unit loses its in-memory state and must not assume a dangerous default.
\end{itemize}

The final project therefore preserves the original \textbf{alarm semantics} while
making the execution \textbf{distributed}, \textbf{role-based}, and
\textbf{failure-aware}. In particular, the implementation directly targets the
assignment requirements of \textbf{at least three nodes}, \textbf{message-only
communication}, and \textbf{safe recovery through}
\texttt{STARTUP\_RECOVERY}.
